% Hierarchical scheduling on a NUMA system: per-CPU runqueues, a load
% balancer, and two memory nodes connected by an interconnect.
% Usage: % Hierarchical scheduling on a NUMA system: per-CPU runqueues, a load
% balancer, and two memory nodes connected by an interconnect.
% Usage: % Hierarchical scheduling on a NUMA system: per-CPU runqueues, a load
% balancer, and two memory nodes connected by an interconnect.
% Usage: % Hierarchical scheduling on a NUMA system: per-CPU runqueues, a load
% balancer, and two memory nodes connected by an interconnect.
% Usage: \input{figures/l07-numa}   (width ~112 mm)
\begin{tikzpicture}[x=1mm, y=1mm, font=\footnotesize\sffamily,
  >={Stealth[length=1.6mm]},
  cpu/.style={draw, fill=ln-accent!22, minimum width=16mm, minimum height=6mm, inner sep=1pt},
  mem/.style={draw, fill=ln-box, minimum width=38mm, minimum height=7mm, inner sep=1pt},
  nd/.style={draw, dashed, rounded corners=2pt, ln-muted},
  lab/.style={font=\scriptsize\sffamily, align=center},
  tk/.style={draw, circle, fill=white, minimum size=2.6mm, inner sep=0pt}]
  % load balancer
  \draw[fill=ln-box] (0,36) rectangle (112,42);
  \node at (56,39) {\iflangja{負荷分散}{load balancing}};
  % per-CPU runqueues and CPUs
  \foreach \x/\n/\k in {13/0/3, 36/1/1, 76/2/2, 99/3/0} {
    \draw[<->] (\x,35.5) -- (\x,29.5);
    \draw[fill=white] (\x-8,24) rectangle (\x+8,29);
    \foreach \i in {1,...,3} \draw (\x+8-4*\i,24) -- ++(0,5);
    \ifnum\k>0 \foreach \i in {1,...,\k} \node[tk] at (\x+10-4*\i,26.5) {}; \fi
    \draw[->] (\x,24) -- (\x,17);
    \node[cpu] (c\n) at (\x,14) {CPU\,\n};
  }
  \node[lab] at (56,26.5) {\iflangja{CPU ごとの\\ランキュー}{per-CPU\\runqueues}};
  % memory nodes
  \node[mem] (m0) at (24.5,-6) {\iflangja{メモリノード 0}{memory node 0}};
  \node[mem] (m1) at (87.5,-6) {\iflangja{メモリノード 1}{memory node 1}};
  \draw[nd] (2,20) rectangle (47,-12);
  \draw[nd] (65,20) rectangle (110,-12);
  \node[lab, anchor=north, text=ln-muted] at (24.5,-12.5) {\iflangja{ノード 0（ソケット 0）}{node 0 (socket 0)}};
  \node[lab, anchor=north, text=ln-muted] at (87.5,-12.5) {\iflangja{ノード 1（ソケット 1）}{node 1 (socket 1)}};
  % local access
  \draw[<->, thick] (c0.south) -- node[lab, right] {\iflangja{ローカル：\\速い}{local:\\fast}} (c0.south |- m0.north);
  \draw (c1.south) -- (c1.south |- m0.north);
  \draw (c2.south) -- (c2.south |- m1.north);
  \draw (c3.south) -- (c3.south |- m1.north);
  % interconnect
  \draw[line width=2pt, ln-muted] (m0.east) -- (m1.west);
  \node[lab, anchor=north, text=ln-muted] at (56,-12.5) {\iflangja{インターコネクト}{interconnect}};
  \draw[ln-muted, densely dotted] (56,-6) -- (56,-12.3);
  % remote access
  \draw[<->, thick, densely dashed, ln-caution] (c2.south west) -- (m0.north east);
  \node[lab, text=ln-caution] at (53.5,10.5) {\iflangja{リモート：\\遅い}{remote:\\slow}};
\end{tikzpicture}
   (width ~112 mm)
\begin{tikzpicture}[x=1mm, y=1mm, font=\footnotesize\sffamily,
  >={Stealth[length=1.6mm]},
  cpu/.style={draw, fill=ln-accent!22, minimum width=16mm, minimum height=6mm, inner sep=1pt},
  mem/.style={draw, fill=ln-box, minimum width=38mm, minimum height=7mm, inner sep=1pt},
  nd/.style={draw, dashed, rounded corners=2pt, ln-muted},
  lab/.style={font=\scriptsize\sffamily, align=center},
  tk/.style={draw, circle, fill=white, minimum size=2.6mm, inner sep=0pt}]
  % load balancer
  \draw[fill=ln-box] (0,36) rectangle (112,42);
  \node at (56,39) {\iflangja{負荷分散}{load balancing}};
  % per-CPU runqueues and CPUs
  \foreach \x/\n/\k in {13/0/3, 36/1/1, 76/2/2, 99/3/0} {
    \draw[<->] (\x,35.5) -- (\x,29.5);
    \draw[fill=white] (\x-8,24) rectangle (\x+8,29);
    \foreach \i in {1,...,3} \draw (\x+8-4*\i,24) -- ++(0,5);
    \ifnum\k>0 \foreach \i in {1,...,\k} \node[tk] at (\x+10-4*\i,26.5) {}; \fi
    \draw[->] (\x,24) -- (\x,17);
    \node[cpu] (c\n) at (\x,14) {CPU\,\n};
  }
  \node[lab] at (56,26.5) {\iflangja{CPU ごとの\\ランキュー}{per-CPU\\runqueues}};
  % memory nodes
  \node[mem] (m0) at (24.5,-6) {\iflangja{メモリノード 0}{memory node 0}};
  \node[mem] (m1) at (87.5,-6) {\iflangja{メモリノード 1}{memory node 1}};
  \draw[nd] (2,20) rectangle (47,-12);
  \draw[nd] (65,20) rectangle (110,-12);
  \node[lab, anchor=north, text=ln-muted] at (24.5,-12.5) {\iflangja{ノード 0（ソケット 0）}{node 0 (socket 0)}};
  \node[lab, anchor=north, text=ln-muted] at (87.5,-12.5) {\iflangja{ノード 1（ソケット 1）}{node 1 (socket 1)}};
  % local access
  \draw[<->, thick] (c0.south) -- node[lab, right] {\iflangja{ローカル：\\速い}{local:\\fast}} (c0.south |- m0.north);
  \draw (c1.south) -- (c1.south |- m0.north);
  \draw (c2.south) -- (c2.south |- m1.north);
  \draw (c3.south) -- (c3.south |- m1.north);
  % interconnect
  \draw[line width=2pt, ln-muted] (m0.east) -- (m1.west);
  \node[lab, anchor=north, text=ln-muted] at (56,-12.5) {\iflangja{インターコネクト}{interconnect}};
  \draw[ln-muted, densely dotted] (56,-6) -- (56,-12.3);
  % remote access
  \draw[<->, thick, densely dashed, ln-caution] (c2.south west) -- (m0.north east);
  \node[lab, text=ln-caution] at (53.5,10.5) {\iflangja{リモート：\\遅い}{remote:\\slow}};
\end{tikzpicture}
   (width ~112 mm)
\begin{tikzpicture}[x=1mm, y=1mm, font=\footnotesize\sffamily,
  >={Stealth[length=1.6mm]},
  cpu/.style={draw, fill=ln-accent!22, minimum width=16mm, minimum height=6mm, inner sep=1pt},
  mem/.style={draw, fill=ln-box, minimum width=38mm, minimum height=7mm, inner sep=1pt},
  nd/.style={draw, dashed, rounded corners=2pt, ln-muted},
  lab/.style={font=\scriptsize\sffamily, align=center},
  tk/.style={draw, circle, fill=white, minimum size=2.6mm, inner sep=0pt}]
  % load balancer
  \draw[fill=ln-box] (0,36) rectangle (112,42);
  \node at (56,39) {\iflangja{負荷分散}{load balancing}};
  % per-CPU runqueues and CPUs
  \foreach \x/\n/\k in {13/0/3, 36/1/1, 76/2/2, 99/3/0} {
    \draw[<->] (\x,35.5) -- (\x,29.5);
    \draw[fill=white] (\x-8,24) rectangle (\x+8,29);
    \foreach \i in {1,...,3} \draw (\x+8-4*\i,24) -- ++(0,5);
    \ifnum\k>0 \foreach \i in {1,...,\k} \node[tk] at (\x+10-4*\i,26.5) {}; \fi
    \draw[->] (\x,24) -- (\x,17);
    \node[cpu] (c\n) at (\x,14) {CPU\,\n};
  }
  \node[lab] at (56,26.5) {\iflangja{CPU ごとの\\ランキュー}{per-CPU\\runqueues}};
  % memory nodes
  \node[mem] (m0) at (24.5,-6) {\iflangja{メモリノード 0}{memory node 0}};
  \node[mem] (m1) at (87.5,-6) {\iflangja{メモリノード 1}{memory node 1}};
  \draw[nd] (2,20) rectangle (47,-12);
  \draw[nd] (65,20) rectangle (110,-12);
  \node[lab, anchor=north, text=ln-muted] at (24.5,-12.5) {\iflangja{ノード 0（ソケット 0）}{node 0 (socket 0)}};
  \node[lab, anchor=north, text=ln-muted] at (87.5,-12.5) {\iflangja{ノード 1（ソケット 1）}{node 1 (socket 1)}};
  % local access
  \draw[<->, thick] (c0.south) -- node[lab, right] {\iflangja{ローカル：\\速い}{local:\\fast}} (c0.south |- m0.north);
  \draw (c1.south) -- (c1.south |- m0.north);
  \draw (c2.south) -- (c2.south |- m1.north);
  \draw (c3.south) -- (c3.south |- m1.north);
  % interconnect
  \draw[line width=2pt, ln-muted] (m0.east) -- (m1.west);
  \node[lab, anchor=north, text=ln-muted] at (56,-12.5) {\iflangja{インターコネクト}{interconnect}};
  \draw[ln-muted, densely dotted] (56,-6) -- (56,-12.3);
  % remote access
  \draw[<->, thick, densely dashed, ln-caution] (c2.south west) -- (m0.north east);
  \node[lab, text=ln-caution] at (53.5,10.5) {\iflangja{リモート：\\遅い}{remote:\\slow}};
\end{tikzpicture}
   (width ~112 mm)
\begin{tikzpicture}[x=1mm, y=1mm, font=\footnotesize\sffamily,
  >={Stealth[length=1.6mm]},
  cpu/.style={draw, fill=ln-accent!22, minimum width=16mm, minimum height=6mm, inner sep=1pt},
  mem/.style={draw, fill=ln-box, minimum width=38mm, minimum height=7mm, inner sep=1pt},
  nd/.style={draw, dashed, rounded corners=2pt, ln-muted},
  lab/.style={font=\scriptsize\sffamily, align=center},
  tk/.style={draw, circle, fill=white, minimum size=2.6mm, inner sep=0pt}]
  % load balancer
  \draw[fill=ln-box] (0,36) rectangle (112,42);
  \node at (56,39) {\iflangja{負荷分散}{load balancing}};
  % per-CPU runqueues and CPUs
  \foreach \x/\n/\k in {13/0/3, 36/1/1, 76/2/2, 99/3/0} {
    \draw[<->] (\x,35.5) -- (\x,29.5);
    \draw[fill=white] (\x-8,24) rectangle (\x+8,29);
    \foreach \i in {1,...,3} \draw (\x+8-4*\i,24) -- ++(0,5);
    \ifnum\k>0 \foreach \i in {1,...,\k} \node[tk] at (\x+10-4*\i,26.5) {}; \fi
    \draw[->] (\x,24) -- (\x,17);
    \node[cpu] (c\n) at (\x,14) {CPU\,\n};
  }
  \node[lab] at (56,26.5) {\iflangja{CPU ごとの\\ランキュー}{per-CPU\\runqueues}};
  % memory nodes
  \node[mem] (m0) at (24.5,-6) {\iflangja{メモリノード 0}{memory node 0}};
  \node[mem] (m1) at (87.5,-6) {\iflangja{メモリノード 1}{memory node 1}};
  \draw[nd] (2,20) rectangle (47,-12);
  \draw[nd] (65,20) rectangle (110,-12);
  \node[lab, anchor=north, text=ln-muted] at (24.5,-12.5) {\iflangja{ノード 0（ソケット 0）}{node 0 (socket 0)}};
  \node[lab, anchor=north, text=ln-muted] at (87.5,-12.5) {\iflangja{ノード 1（ソケット 1）}{node 1 (socket 1)}};
  % local access
  \draw[<->, thick] (c0.south) -- node[lab, right] {\iflangja{ローカル：\\速い}{local:\\fast}} (c0.south |- m0.north);
  \draw (c1.south) -- (c1.south |- m0.north);
  \draw (c2.south) -- (c2.south |- m1.north);
  \draw (c3.south) -- (c3.south |- m1.north);
  % interconnect
  \draw[line width=2pt, ln-muted] (m0.east) -- (m1.west);
  \node[lab, anchor=north, text=ln-muted] at (56,-12.5) {\iflangja{インターコネクト}{interconnect}};
  \draw[ln-muted, densely dotted] (56,-6) -- (56,-12.3);
  % remote access
  \draw[<->, thick, densely dashed, ln-caution] (c2.south west) -- (m0.north east);
  \node[lab, text=ln-caution] at (53.5,10.5) {\iflangja{リモート：\\遅い}{remote:\\slow}};
\end{tikzpicture}
